\documentclass[11pt]{article}
\usepackage[UTF8]{ctex} % 中文支持，建议使用 LuaLaTeX 或 XeLaTeX 编译
\usepackage{graphicx}
\usepackage{amsmath,amssymb}
\usepackage{booktabs}
\usepackage{multirow}
\usepackage{url}

% 参考文献使用 ACL 风格
\bibliographystyle{acl_natbib}

\title{Paper2GraphX：面向科研论文的自动化模型图生成系统}
\author{liuzhou liu}
\date{September 2025}

\begin{document}

\maketitle

% 摘要
\begin{abstract}
本文提出 Paper2GraphX，一个从自然语言 target 到高质量科研模型图的端到端系统。本文档为论文写作框架草案，仅列出章节标题与小节结构。
\end{abstract}

\section{引言}

\section{总体框架}
% 可对应 Paper2Poster.pdf 的整体结构与动机

\section{PaperGraph Schema v0.4}
% 概述语义层、布局层、美学层与辅助层

\section{PaperGraphAgent 架构}
% 本节对应 workflow\_design.md 中的六个 Agent 节点
\subsection{目标解析 Agent（p2g\_target\_analyst）}
\subsection{语义构建 Agent（p2g\_semantic\_constructor）}
\subsection{布局构建 Agent（p2g\_layout\_constructor）}
\subsection{美学设计构建 Agent（p2g\_designer）}
\subsection{布局检查 Agent（p2g\_layout\_checker）}
\subsection{语义检查 Agent（p2g\_semantic\_checker）}

\section{实现与工程细节}
% 关键模块、工具接口、资源管理与可扩展性

\section{实验设置}
% 数据、基线、评价指标与实现细节

\section{实验结果与分析}
% 量化结果、消融研究与可视化案例

\section{相关工作}

\section{讨论与局限}

\section{结论}

\section*{致谢}

\section*{伦理声明}

\bibliography{custom}

\end{document}
